\documentclass[11pt]{article}
\usepackage[T1]{fontenc}
\usepackage{lmodern,amsmath,amssymb,amsthm}
\usepackage[margin=1in]{geometry}
\usepackage[hidelinks]{hyperref}
\newtheorem{theorem}{Theorem}
\newtheorem{lemma}{Lemma}
\newcommand{\C}{\mathbb C}
\newcommand{\N}{\mathbb N}
\newcommand{\PP}{\mathbf P^1(\C)}
\newcommand{\Per}{\operatorname{Per}}
\title{A counterexample to conjecture 00000007767\\[3pt]
\large Squaring and cubing have infinitely many common periodic points}
\author{C0ldSmi1e}
\date{October 5, 2026}
\begin{document}
\maketitle
\begin{abstract}
The complex polynomial maps $f(z)=z^2$ and $g(z)=z^3$ commute and have
infinitely many common periodic points on the complex projective line, but
have no common positive compositional iterate. They therefore refute the
explicit leading equivalence of conjecture 00000007767 in its conventional
complex-dynamical interpretation. The proof uses primitive roots of unity
of orders $5,25,125,\ldots$. The source's unspecified field and iterate-index
conventions are disclosed below; the undefined final uniqueness terminology
is not needed for this disproof.
\end{abstract}

\section{Statement and conventions}
Both the English and Chinese source versions concern commuting polynomials
of degree at least two acting on $\mathbf P^1$. Their first assertion is
\[
 |\Per(f)\cap\Per(g)|=\infty
 \quad\Longleftrightarrow\quad
 \text{$f^{\circ m}=g^{\circ n}$ for some positive integers $m,n$}.
\]
They then add a common-iterating-map assertion and a uniqueness assertion
about a ``symmetrized monomial part'' in a Ritt decomposition. An exact copy
of the bilingual source is bundled as \texttt{conjecture.md}.

The source does not specify a base field. This submission explicitly uses
complex polynomials acting on $\PP=\C\cup\{\infty\}$, the usual classical
complex-dynamical setting. This is a disclosed interpretation, not an
assumption stated verbatim in the source. No conclusion about every possible
arithmetic interpretation is claimed. Positive indices are the standard
nonvacuous meaning of ``common iterate'': allowing $m=n=0$ would give the
identity for every pair. The zeroth iterate is used only to define iteration
recursively. The periodic set means
\[
 \Per(F)=\{P:\text{there exists an integer }r\geq1
                  \text{ with }F^{\circ r}(P)=P\}.
\]
The two return times for a common periodic point may differ. They may also
depend on the point. Periodicity is not replaced by eventual periodicity.
The notation $F^{\circ m}=G^{\circ n}$ means exact equality of self-maps,
not equality up to conjugacy.

We refute the forward implication of the displayed equivalence. The full
source assertion entails that implication, regardless of how its later
uniqueness condition might be made precise. We neither supply a definition
of the undefined final terminology nor claim a replacement Ritt theorem.

\section{The two actual polynomial maps}
Let $f=X^2$ and $g=X^3$ in $\C[X]$. Their projective extensions are
\[
 F([Z:W])=[Z^2:W^2],\qquad G([Z:W])=[Z^3:W^3].
\]
These expressions define maps on projective points: scaling $(Z,W)$ scales
both coordinates by the same nonzero power, and the coordinates cannot
simultaneously vanish. In the affine chart $[z:1]$ the maps are $z\mapsto z^2$
and $z\mapsto z^3$, respectively; $\infty=[1:0]$ is fixed by both.
Their polynomial degrees are two and three. At an affine point both
$F\circ G$ and $G\circ F$ act by $z\mapsto z^6$, and both fix infinity.
Thus the stated degree and commutation hypotheses hold.

For $d\geq1$, let $H_d(z)=z^d$ on the affine chart and $H_d(\infty)=\infty$.
Induction on $r\geq0$ gives
\begin{equation}\label{eq:iterate}
 H_d^{\circ r}(z)=z^{d^r}.
\end{equation}
The induction step is $(z^{d^r})^d=z^{d^{r+1}}$; the case $r=0$ uses
$d^0=1$. This is compositional iteration of the actual maps.

\section{Infinitely many distinct common periodic points}
For each $k\in\N$, put
\[
 N_k=5^{k+1},\qquad \zeta_k=\exp(2\pi i/N_k),\qquad P_k=[\zeta_k:1].
\]
The complex number $\zeta_k$ has exact multiplicative order $N_k$.
Both 2 and 3 are coprime to $N_k$. Euler's theorem gives, for
$t_k=\varphi(N_k)>0$,
\[
 2^{t_k}\equiv1\pmod {N_k},\qquad
 3^{t_k}\equiv1\pmod {N_k}.
\]
Whenever $\zeta^{N}=1$ and $a\equiv1\pmod N$, exponent reduction gives
$\zeta^a=\zeta$. Applying this observation and \eqref{eq:iterate} yields
\[
 F^{\circ t_k}(P_k)=P_k,
 \qquad G^{\circ t_k}(P_k)=P_k.
\]
Every $P_k$ is therefore a common periodic point with a positive return time.

If $P_k=P_\ell$, equality in the affine chart implies
$\zeta_k=\zeta_\ell$. An element has only one exact multiplicative order,
so $5^{k+1}=5^{\ell+1}$. Since powers of 5 are injective in their natural
exponent, $k=\ell$. Hence $k\mapsto P_k$ is an injection from $\N$ into
$\Per(F)\cap\Per(G)$, proving that this one fixed pair of maps has infinitely
many distinct common periodic points. The single point at infinity is not
used to supply infinitude.

\section{No common positive iterate}
Suppose $m,n\geq1$ and $F^{\circ m}=G^{\circ n}$. Evaluating at the affine
point $[2:1]$ and using \eqref{eq:iterate} gives
\[
 2^{2^m}=2^{3^n}.
\]
These are natural numbers viewed in $\C$. Injectivity of the natural-number
embedding into $\C$ transfers the equality back to $\N$. Powers of 2
are injective in their exponent, so $2^m=3^n$. The left side is even because
$m\geq1$, while the right side is odd. This is a contradiction.
This argument excludes every pair of positive indices, not merely a bounded
search or the case $m=n$.

\begin{theorem}
There exist complex polynomials of degree at least two whose projective
extensions commute and have infinitely many common periodic points, but
whose positive compositional iterates are never equal. Consequently the
leading equivalence of conjecture 00000007767 is false in the explicitly
stated complex interpretation.
\end{theorem}
\begin{proof}
Take $f=X^2$ and $g=X^3$. The preceding sections verify all three hypotheses,
the infinitude conclusion, and the failure of every common positive iterate.
\end{proof}

\section{Formalization and reproducibility}
The accompanying Lean 4 project formalizes the complex polynomials and their
extensions to the projective line represented by \texttt{Option Complex}:
\texttt{some z} denotes $[z:1]$ and \texttt{none} denotes $[1:0]$.
Every projective point has exactly one of these forms. The extension of
polynomial evaluation is \texttt{Option.map}; thus the model includes the
point at infinity, actual polynomial evaluation, and compositional iteration.

The namespace is \texttt{Conjecture7767}. The main proof correspondence is:
\begin{center}
\begin{tabular}{ll}
\textbf{Mathematical step} & \textbf{Lean declaration}\\[3pt]
Actual polynomial degrees & \texttt{powerPolynomial\_degree}\\
Composition formula & \texttt{powerMap\_iterate\_affine}\\
Commuting projective maps & \texttt{powerMap\_commute}\\
Positive root-of-unity return time & \texttt{primitiveRoot\_periodic}\\
Distinct roots & \texttt{root\_injective}\\
Infinite intersection & \texttt{common\_periodic\_infinite}\\
No common positive iterate & \texttt{no\_positive\_common\_iterate}\\
All counterexample requirements & \texttt{counterexample}\\
Negation of the leading assertion & \texttt{conjecture\_false}
\end{tabular}
\end{center}

The project pins Lean 4.19.0 and Mathlib v4.19.0. A fresh build, direct
source replay, and the complete \texttt{Check.lean} audit passed with warnings
treated as errors. All 18 authored theorems use only \texttt{propext},
\texttt{Classical.choice}, and \texttt{Quot.sound}. The 11 definitions and
abbreviations were printed and inspected. An exhaustive inspection of all
40 declarations originating in the compiled proof module also passed,
including its 11 generated declarations. There are no custom axioms,
admitted proofs, partial declarations, or authored unsafe definitions. Four
compiler-generated runtime wrappers are explicitly inventoried separately
from the safe logical declarations; none introduces an additional axiom.
Compiler and all nine dependency revisions were checked before and after
the build. The complete logs and hashes accompany the submission.

The auxiliary checker uses exact integer coefficients and rational angles.
It independently compares sparse polynomial iteration with direct integer
iteration, tests all 144 pairs $1\leq m,n\leq12$, and computes every cycle
under doubling and tripling modulo $5,25,125,625$. All 780 root records and
their two complete cycle decompositions are included. For these four orders
the primitive-root counts and periods under either map are respectively
$4,20,100,500$. Together with the root 1, affine zero, and infinity, the
largest tested projective subset has 627 distinct common periodic points.
Normal and optimized Python runs both passed 12,875 explicit checks and
produced byte-identical output. These finite checks establish neither
infinitude nor unbounded nonexistence of common iterates; those statements
are proved in Lean without relying on the Python output.

After installing the pinned Lean toolchain, run the following from the
submission's \texttt{lean} directory (Mathlib cache retrieval requires
network access):
\begin{verbatim}
lake exe cache get
lake build
lake env lean -DwarningAsError=true Conjecture7767.lean
lake env lean -DwarningAsError=true Check.lean
\end{verbatim}
From the \texttt{auxiliary} directory, run \texttt{python3 -B run\_checks.py}.
The full reproduction instructions include the compiled-environment audit
and PDF export. Local validation and the bundled independent semantic review
are distinct from acceptance by a competition maintainer.

\end{document}
